\documentclass[10pt,a4paper]{amsart}
\usepackage[margin=19mm]{geometry}
\usepackage{amsmath,amssymb,mathtools}
\usepackage[scr=rsfs,cal=euler]{mathalfa}
\usepackage{xfrac}
\usepackage[unicode,hidelinks,bookmarks=false]{hyperref}
\newcommand{\ie}{i.e.\ }
\newcommand{\cf}{cf.\ }
\newcommand{\resp}{respectively\ }
\newcommand{\Subsection}{\subsection}
\providecommand{\nobreakdash}{\nobreak\hbox{-}}
\setlength{\emergencystretch}{2em}
\multlinegap0pt
\usepackage{bm}
\usepackage{bbm}
\usepackage{dsfont}
\usepackage{xspace}
\usepackage{cancel}
\usepackage{mathtools}
\usepackage{cleveref}
\usepackage{amsthm}
\makeatletter
\@ifundefined{theorem}{\newtheorem{theorem}{Theorem}}{}
\@ifundefined{lemma}{\newtheorem{lemma}[theorem]{Lemma}}{}
\makeatother
\let\iso=\cong
\let\smash=\wedge

% It is completely ridiculous that this is the way to get reverse ddots
\newcommand{\FF}{\mathbb{F}}
\newcommand{\KQ}{\mathsf{KQ}}
\newcommand{\MZ}{\mathsf{M}\Z}
\newcommand{\R}{\mathbb{R}}
\newcommand{\SH}{\mathbf{SH}}
\newcommand{\Sq}{\operatorname{Sq}}
\newcommand{\Z}{\mathbb{Z}}
\newcommand{\f}{\mathsf{f}}
\newcommand{\pr}{\operatorname{pr}}
\newcommand{\s}{\mathsf{s}}
\newcommand{\vd}{{\tilde{\mathsf{d}}}}
\newcommand{\vs}{\tilde{\mathsf{s}}}
\title[Hermitian $K$-theory and Milnor-Witt motivic cohomology over]{Hermitian $K$-theory and Milnor-Witt motivic cohomology over $\mathbb Z$}
\author{H{\aa}kon Kolderup, Oliver R\"ondigs, Paul Arne {\O}stv{\ae}r}
\date{Source: arXiv:2509.16404v1 (CC BY 4.0)}
\begin{document}
\maketitle

\noindent\emph{Source and licence.} Hermitian $K$-theory and Milnor-Witt motivic cohomology over $\mathbb Z$. Version arXiv:2509.16404v1 (CC BY 4.0), distributed under the Creative Commons Attribution 4.0 International licence, as recorded in the arXiv licence metadata for this exact version and shown on the abstract page https://arxiv.org/abs/2509.16404v1. Full rights notice: licenses/arxiv-2509.16404-CC-BY-4.0.txt.\par\smallskip

\noindent\textit{Editorial scope.} Counted unit: the Lemma whose source label is \texttt{lem:vd-KQ-on-pi} in the article (source0.tex lines 2547--2557, source label \texttt{lem:vd-KQ-on-pi}), delivered together with the article's own complete original proof of it (source0.tex lines 2559--2621, one \texttt{proof} environment of 63 lines). No mathematics is written, repaired or re-derived: statement and proof are the article's own text line for line. Required context retained uncounted: lem:unique-v-operation (lines 2189--2192); lem:hZ (lines 2927--2940). Every definition, display and label of those lines is retained unchanged. Omitted from this package and not redistributed: 1--2188, 2193--2546, 2558--2558, 2622--2926, 2941--3431. Nothing inside a retained excerpt is omitted or altered: every retained body is delivered exactly as the article's lines are written, so no omission receipt of any kind accompanies this package. Citations: the retained excerpts carry no citation command at all, so no bibliography entry of the article is transcribed and no reference is redistributed. Reference adaptation: none was needed and none was made; every reference inside the delivered unit resolves inside it, so no unresolved reference or citation is produced. Printed slips kept literal: no printed word, symbol, hypothesis or number of the counted statement or of its proof was altered. Layout and class: the article's own class is \texttt{article} with packages including geometry, babel, amsmath, amssymb, amsthm, latexsym, stmaryrd, tikz, tikz-cd, array, pdflscape, booktabs, mathtools, tabularx; the wrapper is the lane's pinned \texttt{amsart} editorial wrapper at 10pt on A4, which restates the theorem-like environments the retained text needs and adds the article's own macro definitions that the retained text needs (\texttt{\textbackslash FF}, \texttt{\textbackslash KQ}, \texttt{\textbackslash MZ}, \texttt{\textbackslash R}, \texttt{\textbackslash SH}, \texttt{\textbackslash Sq}, \texttt{\textbackslash Z}, \texttt{\textbackslash f}, \texttt{\textbackslash iso}, \texttt{\textbackslash pr}, \texttt{\textbackslash s}, \texttt{\textbackslash vd}, \texttt{\textbackslash vs}), with extra wrapper packages (bm, bbm, dsfont, xspace, cancel, mathtools, cleveref). The delivered numbering is the unit's own. Assets: no retained span contains a figure environment, a graphics-inclusion command, a PDF-page inclusion, a table or a TikZ picture; no image file is copied into this package, the package declares no asset dependency and no image file is redistributed with it. Licence: version arXiv:2509.16404v1 of this article is distributed under the Creative Commons Attribution 4.0 International licence, as recorded in the arXiv licence metadata for this exact version and shown on the abstract page https://arxiv.org/abs/2509.16404v1. A full rights notice is delivered at licenses/arxiv-2509.16404-CC-BY-4.0.txt. Characters: every retained excerpt is pure ASCII, so the delivered engine, XeLaTeX with its default Latin Modern text font, reports no missing character.
\begin{lemma}\label{lem:unique-v-operation}
In $\SH(\Z)$ there is a unique nonzero map
\[ \vs_0\KQ\to \Sigma\vs_1\KQ \simeq \Sigma^{3,1}\MZ/2\]
\end{lemma}

\begin{lemma}\label{lem:hZ}
The canonical map $h^{\ast,\star}(\Z)\to h^{\ast,\star}(\R)$ is a surjective ring homomorphism whose kernel is concentrated in $\ast=1$ and $\star>1$. 
More precisely, the kernel of $h^{1,w}(\Z)\to h^{1,w}(\R)$ is $\Z/2$ for $w>1$.
If $\tau\epsilon\in h^{1,2}(\Z)$ is the nonzero element in this kernel for $w=2$, 
and $\tau\in h^{0,1}(\Z),\rho\in h^{1,1}(\Z)$ are the usual elements, 
the ring $h^{\ast,\star}(\Z)$ has the following presentation:
\[ 
h^{\ast,\star}(\Z)\iso 
\FF_2[\tau,\rho,\tau\epsilon]/(\tau\epsilon)^2,\tau\epsilon\cdot \rho
\]
Finally, 
$\tau^{2m-1}\cdot \tau\epsilon $ is the image of a generator of the infinite cyclic summand in 
$H^{1,2m+1}(\Z)\cong \Z\oplus \Z/2$.
\end{lemma}

\begin{lemma}\label{lem:vd-KQ-on-pi}
Let $S$ be the spectrum of a field or the ring of integers.
The map 
\[ \pi_{s-(w)}\vd^1_0\KQ \colon \pi_{s-(w)}\vs_{0}\KQ \to \pi_{s-(w)}\Sigma^{3,1}\MZ/2\iso h^{w-s+3,w+1}\]
is zero for all $s<2$ and for all $s\equiv 0,1 \bmod 4$, as well as for all $s$ if $S$ is the spectrum of a field of characteristic $2$.
For $S$ not the spectrum of a field of characteristic $2$, if $s\geq 2$ and $s\equiv 2,3 \bmod 4$, the composition
\[ \pi_{s-(w)}\f_1\vs_0\KQ \to \pi_{s-(w)}\vs_0\KQ \xrightarrow{\pi_{s-(w)}\vd^1_0\KQ} h^{w-s+3,w+1}\]
coincides with the composition
\[ \pi_{s-(w)}\f_1\vs_0\KQ \iso h^{w-s+1,w+1}\xrightarrow{\phi} h^{w-s+3,w+1}\iso \pi_{s-(w)}\Sigma^{3,1}\MZ/2\]
where $\phi(\tau^{s}x)=\tau^{s-2}\rho^2x$ for all $\tau^sx\in h^{w-s+1,w+1}$.
\end{lemma}

\begin{proof}
Note that the target group $\pi_{s-(\star)}\Sigma^{3,1}\MZ/2\cong \mathbf{h}^{\star+1-(s-2),\star+1}$ is zero for $s<2$. 
If $S$ is the spectrum of a field of characteristic $2$, the target of the map in question is zero unless $s=2$. 
In this case the source is the $2$-divisible group $\pi_{2-(w)}\vs_0\KQ\iso H^{w-2,w}$ mapping trivially to the target $h^{w+1,w+1}$.
Hence we may assume $s>1$ and exclude fields of characteristic $2$. 
By the proof of \Cref{lem:unique-v-operation}, the differential $\vd^1_0\KQ$ uniquely corresponds to a map $\f_1\vs_0\KQ\to \Sigma^{3,1}\MZ/2$ factoring as 
\[ \f_1\vs_0\KQ \to \s_1\vs_0(\KQ)\simeq \Sigma^{1,1}\MZ/2\oplus \Sigma^{3,1}\MZ/2 \xrightarrow{\pr} \Sigma^{3,1}\MZ/2\]
where the last map is the projection onto the last summand.
The first slice differential for $\vs_0\KQ$ implies that for $s>1$, $\pi_{s-(w)}\f_1\vs_0\KQ$ is the subgroup
\[ \bigl\lbrace (a,b)\in h^{w-s+1,w+1}\oplus h^{w-s+3,w+1}\colon \Sq^2a=\tau b\bigr\rbrace. \]
If $s\equiv 0,1 \bmod 4$, then $\Sq^2a=0$.
Since multiplication with $\tau$ is always injective (see \Cref{lem:hZ} for the case of $\Z$), the projection above sends this subgroup to zero.
If $s\equiv 2,3(4) \bmod 4$ and $a=\tau^sx$, then $\Sq^2a=\Sq^2(\tau^sx)=\tau^{s-1}\rho^2x$.
The equation $\Sq^2a=\tau b$ then implies $b=\tau^{s-2}\rho^2x$, which gives the result.
%Theorem~\ref{thm:pi-v} shows there is an exact sequence
%\begin{equation}\label{eq:zero-coeff30}
%0 \to \mathbf{h}^{\star-s+1,\star+1} \to 
%\pi_{s-\bideg}\vs_{0}\KQ  \to 
%\mathbf{H}^{\star-s,\star} \to 0
%\end{equation}
%The composite of $\mathbf{h}^{\star-s+1,\star+1}\to\pi_{s-\bideg}\vs_{0}\KQ $ in \eqref{eq:zero-coeff30} and 
%$\pi_{s-(\star)}\vd^1_0\KQ $ yields a $(2,0)$-operation in motivic cohomology with $\mathbb{F}_2$-coefficients, hence it is zero.
%is zero since it results in a natural transformation of 
%motivic cohomology with coefficients in $\mathbb{F}_2$ of weight 0 and dimension two.
%Thus $\pi_{s-(\star)}\vd^1_0\KQ $ factors through the quotient in \eqref{eq:zero-coeff30}, 
%resulting in a $(3,1)$-operation from integer to $\mathbb{F}_2$-coefficients. 
%There is a unique such nontrivial cohomology operation, namely $\Sq^1\Sq^2\pr^\infty_2$.
%The proof of Lemma~\ref{lem:vtomod2weight1} shows that the composite of $\Sq^1\Sq^2\pr^\infty_2$ and the canonical map 
%$\vs_{0}\KQ \to \s_0\vs_{0}\KQ $ (which induces the quotient map in \eqref{eq:zero-coeff30}) is zero as a map of motivic spectra, 
%and not the nonzero map $\vd^1_0$.
%It follows that $\pi_{s-(\star)}\vd^1_0\KQ $ factors through the trivial motivic cohomology operation if 
%$s\equiv 0,3 \bmod 4$, hence is also zero.
%If $s\equiv 1 \bmod 4$, 
%we argue similarly using the exact sequence 
%\begin{equation}\label{eq:zero-coeff1}
%0 \to \mathbf{h}^{\star-s+1,\star+1}/\Sq^2\pr^\infty_2 %\to 
%\pi_{s-\bideg}\vs_{0}\KQ  \to 
%\mathbf{H}^{\star-s,\star} \to 0
%\end{equation}
%in Theorem~\ref{thm:pi-v}.
%In this case, 
%the precomposition of $\pi_{s-(\star)}\vd^1_0\KQ $ with the injection in \eqref{eq:zero-coeff1} yields a 
%motivic cohomology operation after precomposition with the quotient map 
%$\mathbf{h}^{\star-s+1,\star+1}\to \mathbf{h}^{\star-s+1,\star+1}/\Sq^2\pr^\infty_2$. 
%Finally, for $s\equiv 2 \bmod 4$, we use the exact sequence
%\begin{equation}\label{eq:zero-coeff2}
%0 \to \mathbf{h}^{\star-s+1,\star+1} \to 
%\pi_{s-\bideg}\vs_{0}\KQ  \to 
%\ker(\Sq^2\pr^\infty_2\colon \mathbf{H}^{\star-s,\star}\to \mathbf{h}^{\star-s+2,\star+1}) 
%\to 0
%\end{equation}
%from Theorem~\ref{thm:pi-v}.
%Note that $\pi_{s-(\star)}\vd^1_0\KQ $ factors through the quotient in \eqref{eq:zero-coeff2}, 
%which is a submodule of motivic cohomology with integral coefficients.
%Hence, 
%maps from it to $\mathbf{h}^{\star-(s-2)+1,\star+1}$ do not necessarily correspond to 
%motivic cohomology operations.
%In fact, 
%maps from the said quotient to $\mathbf{h}^{\star-(s-2)+1,\star+1}$ correspond to elements in the Ext-group 
%$\Ext^1_{\KMW}(\Sq^2\pr^\infty_2(\mathbf{H}^{\star -s,\star}),\mathbf{h}^{\star-(s-2)+1,\star+1})$ mapping to zero 
%under the canonical map to $\Ext^1_{\KMW}(\mathbf{H}^{\star -s,\star},\mathbf{h}^{\star-(s-2)+1,\star+1})$. 
%\todo{Unfinished} 
\end{proof}

\end{document}
